% Generated draft: requires mathematical and source review.
\documentclass[11pt]{article}
\usepackage[margin=25mm]{geometry}
\usepackage{fontspec}
\setmainfont{Latin Modern Roman}
\usepackage{amsmath,amssymb,mathtools,unicode-math}
\setmathfont{Latin Modern Math}
\usepackage{enumitem,needspace,xurl,hyperref}
\setlength{\parindent}{0pt}
\setlength{\parskip}{5pt}
\setlength{\emergencystretch}{3em}
\newcommand{\N}{\mathbb{N}}
\newcommand{\Z}{\mathbb{Z}}
\newcommand{\Q}{\mathbb{Q}}
\newcommand{\R}{\mathbb{R}}
\newcommand{\C}{\mathbb{C}}
\newcommand{\F}{\mathbb{F}}
\newcommand{\A}{\mathbb{A}}
\newcommand{\PP}{\mathbb{P}}
\newcommand{\E}{\mathbb{E}}
\newcommand{\eps}{\varepsilon}
\newcommand{\dd}{\,\mathrm{d}}
\begin{document}
\section*{UnsolvedMath ID: 20003346}
\textbf{Source problem code:} catalogue entry
\subsection{Definitions and mathematical statement}
This canonical item is the middle of a status sentence after Weiss's Conjecture 48, not the statement of a conjecture. It says that clause (1) was proved in [33] when $\operatorname{rank}_{\Q}G=2$ and in [201] in complete generality, and ends ``We also mention that''. The supplied aid identifies [33] with Chatterjee and Morris, and [201] with Weiss. The missing clause (1) and the unfinished continuation cannot be reconstructed from that sentence alone. In particular, the registry's open label is not evidence that the cited clause remained open.

The supplied companion finite-isogeny comparison may be formulated with its compatibility hypotheses made explicit. Let $\varphi:\widetilde G\to G$ be a finite central $\Q$-isogeny of connected reductive algebraic groups. Put $L_1=\widetilde G(\R)^\circ$ and $L_2=\varphi(L_1)$; using the image avoids assuming surjectivity on all real points. Let compatible arithmetic lattices $\Gamma_1\subset L_1$, $\Gamma_2\subset L_2$ satisfy
\[
 \varphi(\Gamma_1)\subseteq\Gamma_2,\qquad
 [\varphi^{-1}(\Gamma_2)\cap L_1:\Gamma_1]<\infty.
\]
The induced map $F:L_1/\Gamma_1\to L_2/\Gamma_2$ is then a proper finite covering. For a connected real-split torus subgroup $A_1\subset L_1$, set $A_2=\varphi(A_1)$, choose $g_1\in L_1$, $g_2=\varphi(g_1)$, and define the corresponding orbit maps
\[
 o_j:A_j\to L_j/\Gamma_j,\qquad a\longmapsto ag_j\Gamma_j.
\]
Proper means that inverse images of compact sets are compact. Under these hypotheses the aid's comparison gives
\[
 \begin{gathered}
 \operatorname{rank}_{\Q}\widetilde G=\operatorname{rank}_{\Q}G,\qquad
 \operatorname{rank}_{\R}\widetilde G=\operatorname{rank}_{\R}G,\qquad
 \dim A_1=\dim A_2,\\
 o_1\text{ is proper}\quad\Longleftrightarrow\quad o_2\text{ is proper}.
 \end{gathered}
\]
Here split rank is the dimension of a maximal algebraic split torus over the indicated field. Orbit properness concerns maps from the acting groups themselves; changing to an orbit modulo its stabilizer is a different convention and is not silently done here. Finite-index lattice compatibility and the actual real image are essential stated hypotheses.

The aid's worked example is
\[
 q(x)=2x_1x_2+x_3^2-2x_4^2,\qquad
 \operatorname{Spin}(q)\longrightarrow\operatorname{SO}(q).
\]
The hyperbolic plane plus the $\Q$-anisotropic binary form $x_3^2-2x_4^2$ gives $\Q$-rank one, while real signature $(2,2)$ gives real rank two, for both isogenous groups. Thus a strict inequality $\dim A>\operatorname{rank}_{\Q}G$, when used in a rational-rank obstruction, is unchanged by this compatible isogeny comparison; replacing rational rank by real rank would alter it. The absent wording of Conjecture 48(1) is not asserted on the basis of this example.
\subsection{Short English statement}
The record reports proofs of a missing clause of Weiss's conjecture, first in rational rank two and then generally. A supplied companion result says that compatible finite central isogenies preserve rational and real ranks and properness of corresponding split-torus orbit maps. Its Spin-to-SO example has rational rank one but real rank two.
\subsection{Sources}
\begin{itemize}
\item Local UnsolvedMath registry, permanent ID 20003346, exact source code catalogue entry; \texttt{aim-other-notes.json}, notes[152], status commentary following Conjecture 48.
\item American Institute of Mathematics, \emph{Emerging applications of measure rigidity}, \url{https://aimath.org/WWN/measrigid/measrigid.pdf}. Source labels [33] and [201] are retained; full entries are absent. The supplied aid identifies them as Chatterjee--Morris and Weiss, respectively, and attributes Conjecture 48 to Weiss.
\item Supplied research attempt 1: context\_only, problem status not\_a\_problem, rigor label proved, novelty confidence low. It supplies the finite-isogeny audit and quadratic-form example. Machine-generated, not peer reviewed; independently verify before citation. The canonical original remarks and literature fields are empty.
\end{itemize}
\end{document}
